\subsubsection{桥公理升级：\texorpdfstring{$\mathbf{Bridge^{RH}}$}{BridgeRH} 与 RH 的编译}\label{subsec:bridge-rh-axiom-upgrade}
本小节给出 \(\mathbf{Bridge^{RH}}\) 的最小闭合合同：把 RH 目标编译为体系内条件定理，并将“零点事件必须在 \(\mathbf{PW}\) 中可寻址且非 \(\mathsf{NULL}\)”作为显式条款写入桥接接口。

\paragraph{完成化对象}
记完成化函数
\[
\Xi_{\zeta}(s)
:=\frac12\,s(s-1)\,\pi^{-s/2}\Gamma\!\left(\frac{s}{2}\right)\zeta(s),
\]
其零点集恰为黎曼 \(\zeta\) 的非平凡零点（含重数）。RH 可表述为：
\[
\Xi_{\zeta}(s_0)=0,\ 0<\Re(s_0)<1\quad\Longrightarrow\quad \Re(s_0)=\tfrac12.
\]

\begin{remark}[去无穷化等价接口：圆盘零点禁入与 Jensen 缺陷]\label{rem:bridge-rh-disk-jensen}
附录 \ref{app:unit-circle-phase-gate} 给出完整细节；本节仅调用结论接口：通过 Cayley 坐标可把 \(\Xi_\zeta\) 搬运为开圆盘内解析函数 \(F_\zeta\)，并把 RH 判据转写为“盘内无零点”（定理 \ref{thm:app-rh-iff-disk-zero-free}）。等价地，可用 Jensen 缺陷
\[
D_{\zeta}(\varrho)
=\frac{1}{2\pi}\int_{0}^{2\pi}\log|F_{\zeta}(\varrho e^{\mathrm{i}\theta})|\,\dd\theta-\log|F_{\zeta}(0)|
\]
在 \(\varrho\in(0,1)\) 上恒为 \(0\) 来表述（推论 \ref{cor:app-rh-iff-jensen-defect-zero}）。在 Hardy/Smirnov 口径下，内外因子分解 \(F_\zeta=B\cdot S\cdot O\) 将该判据归结为内因子平凡性；这正是桥合同中“隐藏零点数据需外置为可审计轴”的来源。
\end{remark}

\paragraph{当前仍然存在的问题（阻断闭合的关键缺口）}
上述等价式在坐标层面实现了严格去无穷化；然而要据此推出 RH，仍需补齐端点层的刚性控制与桥接闭合。现阶段至少存在如下缺口：
\begin{itemize}
  \item \textbf{(2.1 端点 \(w=-1\) 的微局部刚性缺口)} 圆盘化把 \(t\to\pm\infty\) 合并压到边界端点 \(w=-1\)。当 \(\varrho\uparrow 1\) 时，圆周平均的高虚部采样被迫集中到 \(\theta\approx\pi\) 的窄角邻域（推论 \ref{cor:app-jensen-endpoint-localization}）。因此由 Jensen 缺陷恒零推出 RH 的关键在于：必须证明端点角向微区间内不会产生任何导致 \(D_{\zeta}(\varrho)>0\) 的质量；该环节需要跨尺度、端点局部的统一控制，而非固定半径估计。
  \item \textbf{(2.2 Jensen 缺陷的正项结构导致不可抵消)} Jensen 求和式
  \(D_{\zeta}(\varrho)=\sum_{|a_k|<\varrho}\log(\varrho/|a_k|)\ge 0\)
  具有严格正项结构；任何盘内零点都会在边界层留下不可消去的正下界（命题 \ref{prop:app-jensen-single-zero-lower-bound}）。因此证明目标必须是排除盘内零点数据本身，而非诉诸某种均值抵消机制。
  \item \textbf{(2.3 内外因子分解属于条件口径，且内因子对边界模长不可见)} Hardy/Smirnov 类下的分解 \(F_{\zeta}=B\cdot S\cdot O\) 依赖函数类假设，且边界上几乎处处 \(|BS|=1\) 使内因子数据无法由 \(|F_{\zeta}|\) 直接读取（附录 \S\ref{subsec:unit-disk-inner-outer}）。因此存在两层缺口：其一是证明 \(F_{\zeta}\)（或其归一化版本）确属该类；其二是即便属该类，仍需额外刚性输入以强迫 \(B\equiv 1\)（并进一步排除 \(S\)）。
  \item \textbf{(2.4 尺度链符号已封装，但与缺陷恒零之间缺少闭合推导桥)} 端点精度势 \(y\) 的加法律与尺度转移算子 \(\mathcal{T}\) 的 Mellin 对角化，确实把无限尺度链封装为 \(\zeta\) 的乘子符号；但仍缺少一条闭合论证：将“尺度链的符号结构”转化为“Jensen 缺陷恒零”或“内因子平凡”的刚性结论。具体而言，需要构造一个与 \(D_{\zeta}\) 或 Blaschke 数据同构的可审计对象，使其在 \(\mathcal{T}/\mathcal{K}\) 作用下满足足够强的正性/可逆性/单调性，从而推出盘内零点不可能存在。
  \item \textbf{(2.5 最强去无穷化路线仍是条件性的)} 文稿提出 \(\mathbf{Bridge^{RH}}\) 合同，把 RH 编译为体系内条件定理，并把“每个真实零点 PW-可寻址（不为 \(\mathsf{NULL}\)）”写入桥公理（假设 \ref{assump:bridge-rh}）。在自对偶并进一步假设 \(\Delta(z,u)=\det(I-zB(u))\) 来自有限维矩阵族时，临界线零点可转写为谱穿越判据（命题 \ref{prop:xi-unitary-slice-spectral-crossing}）。该路线在形式上实现极致有限化，但核心困难转移为：如何构造并核验满足桥合同的 \(B(u)\)（以及与 \(\Xi_{\zeta}\) 的严格一致性），从而闭合该编译链路。
\end{itemize}

\begin{remark}[unitary slice 的实化：单位圆判据 \(\Longleftrightarrow\) Cayley 超双曲判据]\label{rem:bridge-rh-cayley-hyperbolicity}
桥合同 \(\mathbf{Bridge^{RH}}\) 的可核验零点读出被类型锁定在 unitary slice（假设 \ref{assump:bridge-rh}-(3)），其几何语言天然落在“单位圆根/单位圆谱”上。
然而这一单位圆几何并非不可约：对任意以单位圆根分布为判据的多项式对象，可用 Cayley 变换把“单位圆全在根”严格转译为“实轴全实根”的一元实\textbf{超双曲性}判据（命题 \ref{prop:xi-rh-cayley-hyperbolicity}），从而允许引入完全实化的证书体系（Sturm 链、Hermite 形式、Hankel 型 PSD 等）。
同时，对应的谱实现也在 Cayley 坐标下合并：酉 Hilbert--P{\'o}lya 与自伴 Hilbert--P{\'o}lya 是同一对象的两种坐标读出（注 \ref{rem:xi-cayley-hilbert-polya-unify}）。
\end{remark}

\begin{remark}[\texorpdfstring{$\mathbf{Bridge^{RH}}$}{BridgeRH} 的有理义务化：完成化迹坐标与幂作用的全有理参数族]\label{rem:bridge-rh-rational-obligationization}
假设 \ref{assump:bridge-rh} 的 (2)--(3) 将“零点事件”的可核验读出类型锁定在 unitary slice，并要求其对幂替换半群动作闭合（与 Chebyshev--Adams 动力学一致）。
为避免与 Hardy 因子分解中的奇异内因子记号混淆，本注中 \(S\) 始终表示完成化迹坐标。
附录 \ref{app:unit-circle-phase-gate} 证明：在 Cayley 紧化坐标 \(x\in\overline{\RR}\) 下，完成化迹坐标具有全有理参数化
\[
S(x)=\iota(x)+\iota(x)^{-1}=\frac{2(1-x^{2})}{1+x^{2}}\in[-2,2],
\]
并且 Chebyshev--Adams 作用 \(S\mapsto C_d(S)\) 存在整系数多项式对 \((A_d,B_d)\in\ZZ[x]^2\) 给出的有理提升 \(\tau_d(x)=B_d(x)/A_d(x)\)，使得严格交换律
\[
C_d\!\bigl(S(x)\bigr)=S\!\bigl(\tau_d(x)\bigr)
\]
成立（命题 \ref{prop:chebyshev-adams-fully-rational}）。
因此，任何仅以 \(S\in[-2,2]\) 为连续实参、并在 \(C_d\) 作用下闭合的 unitary-slice 审计义务，都可等价转写为关于 \(x\in\overline{\RR}\) 的全有理参数族上的义务：其核验输出可被组织为整系数多项式恒等式与有理区间证书。
配合附录 \ref{subsec:unitary-two-integer-certificate} 的二整数 unitary 证书协议，\(\mathbf{Bridge^{RH}}\) 的可验证负担可从“连续三角参数域”降为“有理参数域 + 可回放证书”的机械化核验任务。
\end{remark}

\begin{assumption}[\texorpdfstring{$\mathbf{Bridge^{RH}}$}{BridgeRH}：零点的可寻址桥接合同]\label{assump:bridge-rh}
存在一个由本文 SPG/\(\mathbf{PW}\) 管线诱导的二变量行列式对象与其极限规格，使得：
\begin{enumerate}
  \item \textbf{（目标桥接）}
  存在 \(L>1\) 以及一族（有限层可审计、极限可进入无穷维的）核对象，使其二变量行列式
  \(\Delta(z,u)=\det(I-zB(u))\)（有限维或 Fredholm 口径，定义 \ref{def:fredholm-determinant}）
  满足自对偶 \(\Delta(z,u)=\Delta(uz,1/u)\)，并且其完成化截面
  \(\Xi_{\Delta,L}(s)=\Delta(L^{-s},L^{2s-1})\)
  在整个复平面上与 \(\Xi_{\zeta}(s)\) 一致（或至少具有相同零点集与重数）。
  桥接核验可按命题 \ref{prop:real-arc-sufficiency-unit-disk} 降维：只需在真实协议路径 \(r\uparrow 1\) 的实轴区间 \((0,1)\) 上逐点匹配即可推出单位圆盘内的全盘恒等（从而为极限规格提供解析刚性护栏）。
  \item \textbf{（零点可寻址桥接）}
  对任意满足 \(\Xi_{\zeta}(s_0)=0\) 且 \(0<\Re(s_0)<1\) 的零点 \(s_0\)，存在一个投影词地址 \(W(s_0)\in\mathbf{PW}\)，使其在本文可审计片段内可规约到接口正规形并给出非空值读出（不为 \(\mathsf{NULL}\)，见第 \ref{sec:recursive-addressing} 节），且该读出在闭环接口上验证 \(\Xi_{\Delta,L}(s_0)=0\)（从而把“零点事件”接入地址系统）。
  \item \textbf{（unitary-slice 类型锁定）}
  任意“零点事件”的可核对读出必须走 unitary slice 的审计接口：即其核验可归约为命题 \ref{prop:xi-unitary-slice-spectral-crossing} 的谱穿越口径（\(\theta\)-扫描），从而零点地址的参数必满足 \(|u_L(s_0)|=1\)。
  \item \textbf{（NonCancel/Fredholm-HTF 规格）}
  桥接对象满足注 \ref{rem:positive-fredholm-htf} 的正性 Fredholm-HTF 规格：闭层全纯性、正性/整数重数证书（命题 \ref{prop:resolvent-trace-integer-residue-noncancel}，从而 NonCancel 自动）以及 Abel-first/有限部接口兼容性，从而“零点/极点实体”携带可审计的重数/留数证书并可回归到桥包指纹。
\end{enumerate}
\end{assumption}

\begin{theorem}[\texorpdfstring{$\mathbf{Bridge^{RH}}$}{BridgeRH} \(\Rightarrow\) RH（条件版）]\label{thm:bridge-rh-implies-classical-rh}
在假设 \ref{assump:bridge-rh} 下，RH 成立：任意满足 \(\Xi_{\zeta}(s_0)=0\) 且 \(0<\Re(s_0)<1\) 的零点 \(s_0\) 必满足 \(\Re(s_0)=\tfrac12\)。
\end{theorem}

\begin{proof}
取任意满足 \(\Xi_{\zeta}(s_0)=0\) 且 \(0<\Re(s_0)<1\) 的零点 \(s_0\)。
由假设 \ref{assump:bridge-rh}-(1)，有 \(\Xi_{\Delta,L}(s_0)=0\)。
由假设 \ref{assump:bridge-rh}-(2)，\(s_0\) 在 \(\mathbf{PW}\) 中可寻址且读出非空值。
因此 \(s_0\) 不可能满足 \(\Re(s_0)\neq \tfrac12\)，否则按 \(\mathbf{PW}\) 类型安全定理 \ref{thm:xi-offset-null-type-safety}（等价地也可用假设 \ref{assump:bridge-rh}-(3) 的 \(|u_L(s_0)|=1\) 约束）将导致其在闭包内无地址、读出为 \(\mathsf{NULL}\)，与非空值读出矛盾。
故必有 \(\Re(s_0)=\tfrac12\)。证毕。
\end{proof}

\begin{remark}[仍需闭合的关键缺口：把隐藏盘内数据外置为可审计轴]\label{rem:bridge-rh-open-gaps}
以上编译把“无穷高度/无穷分辨率”的困难压缩为若干明确的接口缺口：它们不再以极限本身出现，而表现为“隐藏数据是否可被强制外置并在闭包内核验”的刚性问题。至少包括：
\begin{enumerate}
  \item \textbf{ACC-Rank 的全域一致性尚未闭合：}需证明最小实现维数 \(d_{\min}^{(\ell)}(\theta)\) 在提升层级与相位加密下对所有 \(\theta\) 统一有界（注 \ref{rem:xi-acc-rank}），而非仅在有限采样或若干实例上可见。
  \item \textbf{分位配额证书与“零点标签”的对象化对齐：}目前 \(\gamma_m(\varepsilon)\) 与 \(B_m(\varepsilon)=\exp(m\gamma_m(\varepsilon))\) 的证书针对折叠纤维尾部（定义 \ref{def:pom-gamma-quantile}），要用于零点事件，需把“零点审计输出”的纤维规模与该证书体系对齐，使 prime-register 配额真正覆盖零点标签的碰撞失败（命题 \ref{prop:xi-completeness-auditable-decomposition} 的第二条约束即以此为输入）。
  \item \textbf{圆盘内因子与 \(\mathbf{PW}\) 可见轴的桥接核验仍是实质性义务：}即使已知 RH 等价于盘内零点数据为空（或 \(B\equiv 1\)），内因子在边界模长上不可见，必须通过正交轴外置；Bridge\(^{RH}\) 合同（假设 \ref{assump:bridge-rh}）把该桥接编译为二变量行列式对象 \(\Delta(z,u)=\det(I-zB(u))\) 的一致性与可寻址条款，其构造与核验仍需闭合。
  \item \textbf{端点放大与“零点信号”的耦合尚待固化：}端点喷射的奇通道线性放大与尺度链对角化机制，为把深层尺度算术化提供了可审计轴（附录 \ref{app:unit-circle-phase-gate} 的端点喷射与 Mellin 对角化段落）；但仍需把“深层盘内零点数据”在该放大下转写为可被 unitary-slice/\(\mathbf{PW}\)/ACC-Rank 捕获的稳定证书事件，而非仅得到分辨率处理的算术化外壳。
  \item \textbf{寄存器最小性与径向模式排除仍需严格化：}推论 \ref{cor:fold-tail-patch-incomplete} 与定理 \ref{thm:fold-gauge-anomaly-max} 已给出“局部末端补丁不可能补足跨分辨率一致性”的硬下界，并在命题 \ref{prop:xi-null-three-way}--注 \ref{rem:xi-register-minimality-no-tail-patch} 中将其解释为记录轴的非局部强迫；要将其提升为“\(\mathbf{PW}\) 闭包内不允许出现与 \(\theta\) 解耦的连续径向寄存器”的严格定理，还需把 \(\mathbf{PW}\) 的资源预序/代价函子与该全域下界在同一逻辑层面拼接。
\end{enumerate}
\end{remark}
